\subsection{Software-Architektur und verwendete Werkzeuge}

Die automatisierte Datenverarbeitung und Visualisierung erfolgt in der Programmiersprache \textit{Python (Version 3.11.15)} \cite{zotero-item-49}. Zur Bereitstellung einer reproduzierbaren und systemunabhängigen Entwicklungsumgebung kommt die Distributionsplattform \textit{Anaconda (Version 26.1.1)} \cite{zotero-item-50} zum Einsatz. Somit können Versionskonflikte zwischen den teils komplexen C-basierten Abhängigkeiten der verwendeten Geodaten-Bibliotheken vermieden werden. Die verwendeten Bibliotheken sowie ihr Einsatzgebiet sind in Tabelle \ref{tab:python-lib} dargestellt.



\begin{table}[htbp]
\centering
\caption{Übersicht der eingesetzten Softwarekomponenten und Python-Bibliotheken}
\label{tab:python-lib}
\small
\begin{tabular*}{\textwidth}{@{\extracolsep{\fill}} l p{2.2cm} p{1.2cm} p{4.2cm} p{4.0cm} @{}}
\toprule
\textbf{Kategorie} & \textbf{Bibliothek} & \textbf{Version} & \textbf{Funktion} & \textbf{Referenz} \\
\midrule
\addlinespace
\textbf{Geodaten} & GeoPandas \newline Shapely \newline Rasterio & 1.1.3 \newline 2.1.2 \newline 1.4.4 & Vektor- und Rasterdatenverarbeitung; räumliche Referenzierung. & \cite{zotero-item-51} \newline \cite{zotero-item-52} \newline \cite{zotero-item-53} \\
\addlinespace
\textbf{Punktwolken} & PDAL \newline Laspy & 2.10.0 \newline 2.7.0 & Verarbeitung, Filterung und Einlesen der originalen LiDAR-Daten & \cite{zotero-item-54} \newline \cite{zotero-item-55} \\
\addlinespace
\textbf{Segmentierung} & SciPy \newline scikit-image & 1.14.0 \newline 0.26.0 & Mathematische Algorithmen und Raster-Segmentierung & \cite{zotero-item-70} \newline \cite{zotero-item-72} \\
\addlinespace
\textbf{Visualisierung} & PyVista & 0.48.4 & Interaktives 3D-Rendering und visuelle Qualitätskontrolle & \cite{zotero-item-56}  \\
\bottomrule
\end{tabular*}
\end{table}

Abseits der Codeumgebung war eine visuelle Qualitätskontrolle der räumlichen Daten unerlässlich. Hierfür wurde die Open-Source-Software \textit{CloudCompare (Version 2.13.2)} \cite{zotero-item-57} eingesetzt, die ein performantes Rendering der ALS-Datensätze ermöglicht.